\begin{abstract}
We prove that learned semantic representations alone cannot achieve zero-error identity recovery. When distinct entities map to the same representation, no decoder, regardless of architecture, training, or compute, can distinguish them without additional bits. If $A_\pi$ is the maximum representation-collision multiplicity, then the minimum worst-case auxiliary description is exactly $\log_2 A_\pi$ bits, and this bound is tight.

The coding laws are exact. For repeated independent uses of the same representation map, the minimum block budget is $L_t^\star=t\log_2 A_\pi$. In the adaptive regime, the optimal fiber budget is $\lceil \log_2 |\pi^{-1}(u)| \rceil$, any feasible adaptive scheme has expected length at least the expectation of that quantity and hence at least $H(C\mid U)/\log 2$, and standard fiberwise prefix coding achieves expected length at most $H(C\mid U)/\log 2 + 1$. Under a uniform source on a collision block of size $a$, subcritical fixed-length budgets satisfy the explicit error floor $D \ge 1-2^L/a$.

The same ambiguity-fiber geometry yields exact task-sufficiency, helper-view, and factorized-representation consequences. We instantiate these generic theorems in a concrete formal identity domain, proving that symbolic handles are information-theoretically necessary, not optional, for zero-error open-world identity recovery (\LH{NSL2}, \LH{NSL3}). The instantiated domain is a witness, not a restriction: the necessity claim applies to any learned representation with collisions. The pair consisting of a neural encoder and an injective symbolic handle achieves zero-error identity recovery if and only if the handle distinguishes entities within each semantic fiber (\LH{NSL1}--\LH{NSL6}). The identity bits must be paid somewhere, either in the encoder or in the symbolic handle. No free lunch.

\textbf{Keywords:} semantics-aware compression, zero-error coding, neurosymbolic systems, learned representations, side information
\end{abstract}
